\section{Related work}
\label{sec:related-work}

\nospacestitle{Rollout acceleration with algorithmic modifications. \,}
Several works propose algorithmic changes to accelerate rollout,
e.g., skipping the tail generations~\cite{DBLP:journals/corr/abs-2501-12599,DBLP:journals/corr/abs-2504-13914,rhymerl}.
Such modifications essentially make the training 
a variant of off-policy (vs. on-policy of our targeted setup),
which has unstable convergence in theory~\cite{DBLP:journals/corr/abs-2506-20520,DBLP:journals/corr/abs-2503-14286,DBLP:journals/corr/abs-2405-08448}
and {we observe that they are not always adopted in production~\cite{seer,deepseek-r1}.}
%{\sys} is orthogonal to these methods:
%they do not change the LLM generation process, whereas we directly accelerate
%it without requiring additional GPUs.

\stitle{Rollout acceleration for resource-restrict setups. \,}
Recent works like Seer~\cite{seer} 
and RollPacker~\cite{rollpacker} address long-tail generation 
in resource-constraint setups,
i.e., each worker has limited GPU capacity to hold the 
assigned batch. 
% \remove{For example, each worker in Seer~\cite{seer} can execute a large batch size of {600--800} (in its Table 5).}
As a result, the batch must be further divided into sub-batches where 
each sub-batch is executed sequentially,
and their optimizations lie in reordering or repacking
sub-patches to reduce GPU wastes between sub-batch executions. 
In contrast, {\sys} accelerates long-tail generations where the batch can
fit in the GPU memory of each worker, a common production setup (see {\fig{fig:motiv-batchsize}} (a)).
In general, executing one large batch is faster than sequentially
executing the divided sub-batch (if with sufficient resources) 
because the generation time grows sub-linearly with the batch size 
(see {\fig{fig:motiv-batchsize}} (b)), 
not to mention the possible switching overhead between sub-batches 
with packing techniques~\cite{rollpacker}.

\stitle{Concurrent works on accelerating rollout via speculative decoding. \,}
A few concurrent works 
share a similar high-level idea of using speculative decoding to accelerate rollout~\cite{tlt,seer}.
{\sys} further addresses the low efficiency issue in batch configurations that are unfriendly
to speculative decoding,
and the unknown best drafter selection issue across requests.


The closest work to ours is TLT~\cite{tlt}, 
which selects training-based drafter (EAGLE~\cite{eagle,eagle-2,eagle-3}) 
for speculative rollout. 
It does not consider the large batch efficiency issue and 
its performance can be further improved with our decoupled speculation.
Moreover, our draft ladder with Fastest-of-N speculation 
further provides two advantages over its chosen EAGLE:
(1) EAGLE is more challenging to deploy than the model- and n-gram-based methods 
in our ladder because it requires extra training to match the target model~\cite{eagle-2,eagle-3}, 
which is not always available for the trained model or requires extra tuning; and
(2) we found it has a lower acceptance length than the drafters 
we considered using the models provided by TLT without prompt tuning for EAGLE~\cite{tlt-code} (see \fig{fig:motiv-acc-len}).
Finally, TLT adopts approximate sampling to improve the acceptance rate 
at the cost of training accuracy~\cite{DBLP:journals/corr/abs-2510-26475}.
In contrast, {\sys} uses exact matching to ensure the correctness of the rollout responses.

\stitle{Speculative decoding in inference. \,}
The optimization goal of using speculative decoding 
in rollout is fundamentally different from that in inference systems
~\cite{spd,specinfer,pearl,swiftspec,DBLP:conf/icml/LiuHBCDS024,turbospec}. 
Specifically, rollout prioritizes overall batch execution efficiency, 
where individual request latencies are less critical, 
while inference systems must ensure low latencies for all requests. 
This introduces new challenges, which we address with decoupled and 
Fastest-of-N speculation.

